% Л5 -- Коллекции, IEnumerable и LINQ. Slides.
%
% Companion to the spoken script in ../L05-collections-and-linq.md, which is
% the source of truth: the slides carry diagrams, code and short claims, and
% the delivery lives in the script only -- this deck carries no speaker
% notes. Slide text is in Russian because that is the content; everything
% else in this file is English.
%
% Build:  make          (latexmk -pdf; engine: pdflatex)
%
% Every listing was compiled and run on dotnet SDK 10.0.112 in a stock
% `dotnet new console`; outputs shown or described are real. See
% ../CLAUDE.md for the verification table.
%
% House rules: ../CLAUDE.md and ../../../CLAUDE.md. Theme: ../../../../theme/.

\documentclass[aspectratio=169,10pt,t]{beamer}

\usetheme{dotnet}
\usetikzlibrary{arrows.meta,positioning,calc}
\usepackage{array}   % ragged-right table columns

\dotnetlecture{Л5}

\title[Л5. Коллекции, IEnumerable, LINQ]{Коллекции, IEnumerable и LINQ}
\author{Дмитрий Курочкин}
\date{}

% A listing set one step under \scriptsize, for two columns of code side by
% side. Smaller than this does not carry to the back of the room.
\newcommand{\twocol}{\ttfamily\fontsize{7.0}{8.2}\selectfont}

\begin{document}

% ------------------------------------------------------------------ title --
\begin{frame}[plain]
  \titlepagednc{Лекция 5}{Коллекции, IEnumerable и LINQ}{}{Дмитрий
    Курочкин}
\end{frame}

% ===================================================== 1. DATA STRUCTURES ==
\section{1. Данные в памяти}

\begin{frame}[plain]
  \sectionopener{1}{Данные в памяти}{Сначала структуры без .NET, потом те
    же имена в библиотеке.}
\end{frame}

\begin{frame}{Массив}
  \slidetop
  \claim{Непрерывный блок памяти: адрес \hi{= начало + индекс × размер}.}

  \begin{center}
    \begin{tikzpicture}[font=\small,
      cell/.style={draw=hairline,line width=0.8pt,fill=codebg,
                   minimum width=1.4cm,minimum height=0.9cm,inner sep=0pt}]
      \foreach \v [count=\i from 0] in {10,20,30,40,50,60} {
        \node[cell] (c\i) at (\i*1.4,0) {\texttt{\v}};
        \node[muted,font=\scriptsize\ttfamily,above=0.25em] at (c\i.north)
          {\i};
      }
      \draw[{Straight Barb[length=1.8mm]}-{Straight Barb[length=1.8mm]},
        accent,line width=0.6pt] ([yshift=-0.55cm]c0.south) --
        node[below,font=\scriptsize,text=accent]{длина фиксирована}
        ([yshift=-0.55cm]c5.south);
    \end{tikzpicture}
  \end{center}

  \slidegap
  \aside{Доступ по номеру --- одна операция, независимо от размера. Вставка
    и удаление в середине --- сдвиг \lo{всех} элементов после точки.}
\end{frame}

\begin{frame}{Динамический массив}
  \slidetop
  \claim{Два числа вместо одного: \hi{Count} занято, \hi{Capacity}
    выделено.}

  \begin{center}
    \begin{tikzpicture}[font=\small,
      used/.style={draw=accent,line width=1.0pt,fill=codebg,
                   minimum width=1.3cm,minimum height=0.9cm,inner sep=0pt},
      free/.style={draw=hairline,dashed,line width=0.8pt,fill=paper,
                   minimum width=1.3cm,minimum height=0.9cm,inner sep=0pt}]
      \foreach \i in {0,1,2,3} \node[used] (u\i) at (\i*1.3,0) {};
      \foreach \i in {4,5,6,7} \node[free] (f\i) at (\i*1.3,0) {};
      \draw[accent,line width=0.6pt,{Bar[width=1.6mm]}-{Bar[width=1.6mm]}]
        ([yshift=0.55cm]u0.north) -- node[above,font=\scriptsize,text=accent]
        {Count = 4} ([yshift=0.55cm]u3.north);
      \draw[muted,line width=0.6pt,{Bar[width=1.6mm]}-{Bar[width=1.6mm]}]
        ([yshift=-0.55cm]u0.south) -- node[below,font=\scriptsize,text=muted]
        {Capacity = 8} ([yshift=-0.55cm]f7.south);
    \end{tikzpicture}
  \end{center}

  \slidegap
  \aside{Пока \texttt{Count < Capacity} --- добавление это запись в
    следующую ячейку. Capacity кончилась --- новый массив вдвое больше,
    копия, и старый в мусор. Это \hi{амортизированная} сложность, не
    постоянная.}
\end{frame}

\begin{frame}{Массив vs динамический массив}
  \slidetop
  {\renewcommand{\arraystretch}{1.6}
  \begin{tabular}{@{}p{4.4cm}>{\centering\arraybackslash}p{3.6cm}%
    >{\centering\arraybackslash}p{3.6cm}@{}}
    & \colhead{alt}{массив} & \colhead{accent}{дин. массив} \\
    доступ по индексу    & \hi{O(1)} & \hi{O(1)} \\
    добавление в конец   & \lo{недоступно} & \hi{O(1)} амортиз. \\
    вставка в середину   & $O(n)$ сдвиг    & $O(n)$ сдвиг \\
    память                & без запаса      & с запасом \\
  \end{tabular}}

  \slidegap
  \aside{Динамический массив снимает фиксированный размер, но не снимает
    цену вставки не в конец.}
\end{frame}

\begin{frame}{Связный список}
  \slidetop
  \claim{Вставка и удаление рядом с \hi{известным} узлом --- одна операция;
    доступа по номеру нет.}

  \begin{center}
    \begin{tikzpicture}[font=\footnotesize\ttfamily,
      nd/.style={draw=hairline,line width=0.8pt,fill=codebg,
                 minimum width=1.1cm,minimum height=0.8cm,inner sep=0pt},
      ar/.style={-{Straight Barb[length=1.8mm]},muted,line width=0.6pt}]
      \node[nd] (a) at (0,0)      {10};
      \node[nd] (b) at (3.4,0.6)  {20};
      \node[nd] (c) at (6.6,-0.4) {30};
      \node[nd] (d) at (10.0,0.3) {40};
      \draw[ar] (a) -- (b);
      \draw[ar] (b) -- (c);
      \draw[ar] (c) -- (d);
    \end{tikzpicture}
  \end{center}

  \slidegap
  \aside{Узлы разбросаны по куче как угодно. Чтобы попасть к элементу,
    нужен проход по ссылкам от начала --- это \lo{линейно}.}
\end{frame}

\begin{frame}{Один указатель или два}
  \slidetop
  \begin{center}
    \begin{tikzpicture}[font=\footnotesize,
      nd/.style={draw=hairline,line width=0.8pt,fill=codebg,
                 minimum width=1.6cm,minimum height=0.9cm,inner sep=0pt},
      ptr/.style={draw=muted,circle,fill=paper,minimum size=0.42cm,
                  inner sep=0pt,font=\tiny}]
      \node[nd] (a) at (0,0) {значение};
      \node[ptr] (pa) at ([xshift=0.55cm,yshift=0.3cm]a.east) {N};
      \node[muted,font=\scriptsize,below=0.5em] at (a.south)
        {однонаправленный};

      \node[nd] (b) at (6.0,0) {значение};
      \node[ptr] (pb1) at ([xshift=0.55cm,yshift=0.3cm]b.east) {N};
      \node[ptr] (pb2) at ([xshift=0.55cm,yshift=-0.3cm]b.east) {P};
      \node[muted,font=\scriptsize,below=0.5em] at (b.south)
        {двунаправленный};
    \end{tikzpicture}
  \end{center}

  \slidegap
  \aside{Лишний кружок --- лишний указатель и лишняя память на \hi{каждый}
    узел. В обмен: удаление по одной ссылке на сам узел, без поиска
    \lo{соседей}.}
\end{frame}

\begin{frame}{Хеш-таблица}
  \slidetop
  \claim{Адрес не по номеру, а по значению ключа: \hi{хеш → номер
    корзины}.}

  \begin{center}
    \begin{tikzpicture}[font=\small,
      cell/.style={draw=hairline,line width=0.8pt,fill=codebg,
                   minimum width=1.3cm,minimum height=0.8cm,inner sep=0pt},
      chain/.style={draw=hairline,line width=0.7pt,fill=paper,
                    minimum width=1.0cm,minimum height=0.6cm,inner sep=0pt},
      ar/.style={-{Straight Barb[length=1.8mm]},accent,line width=0.6pt}]
      \foreach \i in {0,1,2,3,4} \node[cell] (c\i) at (\i*1.3,0) {};
      \node[accent,font=\scriptsize\ttfamily,above=0.9cm] (k) at (c2.north)
        {ключ};
      \draw[ar] (k) -- node[right,font=\scriptsize,text=accent]{хеш} (c2.north);
      \node[chain] (h1) at ([yshift=-0.8cm]c2.south) {};
      \draw[muted,line width=0.6pt] (c2.south) -- (h1.north);
      \node[alt,font=\scriptsize,anchor=west] at ([xshift=0.6cm]h1.east)
        {коллизия: \texttt{Equals} в цепочке};
    \end{tikzpicture}
  \end{center}

  \slidegap
  \aside{Хорошая хеш-функция --- короткие цепочки. Плохая --- всё в одной
    корзине, линейный перебор.}
\end{frame}

\begin{frame}{Бинарное дерево поиска}
  \slidetop
  \claim{Слева меньше, справа больше: поиск --- это \hi{бинарный поиск} по
    ссылкам.}

  \begin{center}
    \begin{tikzpicture}[font=\small\ttfamily,level distance=1.1cm,
      sibling distance=2.6cm,
      nd/.style={draw=hairline,line width=0.8pt,fill=codebg,circle,
                 minimum size=0.9cm,inner sep=0pt,font=\small\ttfamily}]
      \node[nd] {30}
        child { node[nd] {15}
          child { node[nd] {5} }
          child { node[nd] {20} } }
        child { node[nd] {45}
          child { node[nd] {35} }
          child { node[nd] {60} } };
    \end{tikzpicture}
  \end{center}

  \slidegap
  \aside{Сбалансированное --- глубина $\approx \log n$, поиск и вставка за
    \hi{логарифм}. Собранное неудачно --- вырождается в связный список.}
\end{frame}

\begin{frame}{Пять структур рядом}
  \slidetop
  {\renewcommand{\arraystretch}{1.5}\footnotesize
  \begin{tabular}{@{}p{2.6cm}>{\centering\arraybackslash}p{1.8cm}%
    >{\centering\arraybackslash}p{1.8cm}>{\centering\arraybackslash}p{1.8cm}%
    >{\centering\arraybackslash}p{1.8cm}>{\centering\arraybackslash}p{1.8cm}@{}}
    & массив & дин.\ массив & связ.\ список & хеш-табл. & дерево \\[0.3em]
    по номеру   & \hi{O(1)} & \hi{O(1)} & $O(n)$    & --- & ---\\
    по значению & $O(n)$    & $O(n)$    & $O(n)$    & \hi{O(1)} & $O(\log n)$\\
    вставка     & $O(n)$    & \hi{O(1)}$^*$ & \hi{O(1)}$^\dagger$ & \hi{O(1)} & $O(\log n)$\\
    порядок     & да        & да        & да        & \lo{нет}  & да \\
    память      & плотная   & с запасом & служебная & служебная & служебная \\
  \end{tabular}\par
  \vspace{0.6em}
  {\scriptsize\color{muted}$^*$в конец, амортизированно\quad
    $^\dagger$у известного узла}}

  \slidegap
  \aside{Каждая следующая платит за то, чего не было у предыдущей.}
\end{frame}

\begin{frame}{Стек и очередь}
  \slidetop
  \claim{Это не структура данных, а \hi{правило доступа} к ней.}

  \begin{center}
    \begin{tikzpicture}[font=\small,
      cell/.style={draw=hairline,line width=0.8pt,fill=codebg,
                   minimum width=1.1cm,minimum height=0.8cm,inner sep=0pt},
      ar/.style={-{Straight Barb[length=1.8mm]},line width=0.7pt}]
      \foreach \i in {0,1,2} \node[cell] (s\i) at (\i*1.1,0) {};
      \draw[ar,accent] ([yshift=0.5cm]s2.north) -- (s2.north);
      \draw[ar,accent] (s2.north) -- ([yshift=0.5cm]s2.north);
      \node[accent,font=\scriptsize,above=0.9cm] at (s2.north)
        {LIFO: один конец};
      \node[muted,font=\scriptsize,below=0.5em] at (s1.south) {стек};

      \foreach \i in {0,1,2} \node[cell] (q\i) at ([xshift=5.4cm]\i*1.1,0) {};
      \draw[ar,alt] ([yshift=0.5cm]q0.north) -- (q0.north);
      \draw[ar,alt] (q2.north) -- ([yshift=0.5cm]q2.north);
      \node[alt,font=\scriptsize,above=0.9cm] at ([xshift=5.4cm]1.1,0.4)
        {FIFO: разные концы};
      \node[muted,font=\scriptsize,below=0.5em] at ([xshift=5.4cm]q1.south)
        {очередь};
    \end{tikzpicture}
  \end{center}

  \slidegap
  \aside{Ни то, ни другое не говорит, \lo{как} элементы хранятся внутри ---
    только в каком порядке их отдают.}
\end{frame}

\begin{frame}{Один и тот же стек на разных структурах}
  \slidetop
  \versus{на массиве}{%
    \begin{cbul}{accent}
      \item добавить/убрать с конца
      \item массив это умеет без сдвига
    \end{cbul}}{на связном списке}{%
    \begin{cbul}{alt}
      \item добавить/убрать с начала
      \item одна операция по ссылке
    \end{cbul}}

  \slidegap
  \hrulethin{0.42}
  \rulegap
  {\large Стек и очередь можно построить на \hi{любой} структуре данных с
    прошлых слайдов.}
\end{frame}

\begin{frame}{Массив → \texttt{T[]}}
  \slidetop
  \claim{Единственная структура с собственными инструкциями в \hi{IL}.}

  \begin{center}
  \begin{tikzpicture}
    \node[draw=hairline,line width=0.8pt,fill=codebg,inner sep=0.6em] {%
      \begin{tabular}{@{}ll@{}}
        \texttt{\hi{newarr}}  & создать массив \\
        \texttt{\hi{ldlen}}   & узнать длину \\
        \texttt{\hi{ldelem}}  & взять элемент \\
        \texttt{\hi{stelem}}  & положить элемент \\
      \end{tabular}};
  \end{tikzpicture}
  \end{center}

  \slidegap
  \aside{Всё остальное в библиотеке --- список, словарь, очередь ---
    построено поверх массива обычным кодом на C\#.}
\end{frame}

\begin{frame}{Один индекс или два}
  \slidetop
  \begin{minipage}[t]{0.47\textwidth}
    \colhead{accent}{int[]}\par\vspace{0.6em}
    \begin{tikzpicture}[font=\small,
      cell/.style={draw=hairline,line width=0.7pt,fill=codebg,
                   minimum width=0.85cm,minimum height=0.75cm,inner sep=0pt}]
      \foreach \i in {0,...,4} \node[cell] (c\i) at (\i*0.85,0) {};
    \end{tikzpicture}\par\vspace{0.5em}
    \aside{адрес = начало + \texttt{i} × размер}
  \end{minipage}\hfill
  \begin{minipage}[t]{0.47\textwidth}
    \colhead{alt}{int[,]}\par\vspace{0.6em}
    \begin{tikzpicture}[font=\small,
      cell/.style={draw=hairline,line width=0.7pt,fill=codebg,
                   minimum width=0.85cm,minimum height=0.6cm,inner sep=0pt}]
      \foreach \r in {0,...,2}
        \foreach \c in {0,...,2}
          \node[cell] at (\c*0.85,-\r*0.6) {};
    \end{tikzpicture}\par\vspace{0.5em}
    \aside{адрес = начало + (\texttt{row}×\texttt{width}+\texttt{col}) ×
      размер}
  \end{minipage}

  \slidegap
  \aside{Один непрерывный блок в обоих случаях --- рантайм просто считает
    адрес по \hi{двум} индексам вместо одного.}
\end{frame}

\begin{frame}[fragile]{Динамический массив → \texttt{List<T>}}
  \slidetop
  \claim{Та же пара чисел, теперь с именами: \hi{Count} и \hi{Capacity}.}

\begin{lstlisting}[style=term,basicstyle=\ttfamily\small]
 1: capacity 4      5: capacity 8      9: capacity 16
 4: capacity 4      8: capacity 8     16: capacity 16
                                      17: capacity 32
\end{lstlisting}

  \slidegap
  \aside{Удвоение на каждом переполнении. Знаете размер заранее ---
    передайте его в конструктор, и копирований не будет вообще.}
\end{frame}

\begin{frame}{Связный список → \texttt{LinkedList<T>}}
  \slidetop
  \claim{Ровно \hi{двунаправленный} вид --- другого в библиотеке нет.}

  \begin{center}
    \begin{tikzpicture}[font=\footnotesize\ttfamily,
      nd/.style={draw=hairline,line width=0.8pt,fill=codebg,
                 minimum width=1.3cm,minimum height=0.8cm,inner sep=0pt},
      ar/.style={-{Straight Barb[length=1.6mm]},muted,line width=0.6pt}]
      \node[nd] (a) at (0,0)   {10};
      \node[nd] (b) at (3.2,0) {20};
      \node[nd] (c) at (6.4,0) {30};
      \draw[ar] ([yshift=0.15cm]a.east) -- ([yshift=0.15cm]b.west);
      \draw[ar] ([yshift=-0.15cm]b.west) -- ([yshift=-0.15cm]a.east);
      \draw[ar] ([yshift=0.15cm]b.east) -- ([yshift=0.15cm]c.west);
      \draw[ar] ([yshift=-0.15cm]c.west) -- ([yshift=-0.15cm]b.east);
    \end{tikzpicture}
  \end{center}

  \slidegap
  \aside{Вставка и удаление у \hi{известного} узла --- одна операция. Узел
    нужно найти --- это уже линейно.}
\end{frame}

\begin{frame}[plain]
  \vote{Список на массиве или связный список --- что быстрее сложить
    миллион чисел?}
\end{frame}

\begin{frame}{Миллион элементов: одна сложность, разное время}
  \slidetop
  \begin{center}
    \begin{tikzpicture}[font=\scriptsize]
      \node[draw=accent,line width=1.2pt,fill=codebg,minimum width=8.6cm,
        minimum height=0.8cm] (arr) at (0,1.3) {};
      \node[accent,font=\scriptsize,below=0.3em] at (arr.south)
        {\texttt{List<T>}: лента, в кэше};

      \foreach \x/\y in {-4.0/-0.5,-3.0/-0.1,-2.2/-0.7,-1.0/-0.2,-0.2/-0.8,
                          0.9/-0.3,1.7/-0.9,2.8/-0.4,3.6/-1.0,4.3/-0.2} {
        \node[draw=alt,line width=0.8pt,fill=codebg,minimum width=0.5cm,
          minimum height=0.5cm] at (\x,\y) {};
      }
      \node[alt,font=\scriptsize] at (0,-1.5)
        {\texttt{LinkedList<T>}: острова в куче};
    \end{tikzpicture}
  \end{center}

  \slidegap
  \aside{Оба \hi{$O(n)$}. Массив читается блоками процессора; переход по
    ссылке --- чтение из случайного места памяти, мимо кэша.}
\end{frame}

\begin{frame}{Хеш-таблица → \texttt{Dictionary<TKey, TValue>}}
  \slidetop
  \begin{center}
    \begin{tikzpicture}[font=\small,
      cell/.style={draw=hairline,line width=0.8pt,fill=codebg,
                   minimum width=1.3cm,minimum height=0.8cm,inner sep=0pt},
      chain/.style={draw=hairline,line width=0.7pt,fill=paper,
                    minimum width=1.0cm,minimum height=0.6cm,inner sep=0pt},
      ar/.style={-{Straight Barb[length=1.8mm]},accent,line width=0.6pt}]
      \foreach \i in {0,1,2,3,4} \node[cell] (c\i) at (\i*1.3,0) {};
      \node[accent,font=\scriptsize\ttfamily,above=1.1cm] (k) at (c2.north)
        {key};
      \draw[ar] (k) -- node[right,font=\scriptsize,text=accent]
        {\texttt{GetHashCode}} (c2.north);
      \node[chain] (h1) at ([yshift=-0.9cm]c2.south) {};
      \draw[muted,line width=0.6pt] (c2.south) -- (h1.north);
      \node[alt,font=\scriptsize,anchor=west] at ([xshift=0.6cm]h1.east)
        {\texttt{Equals} сравнивает внутри корзины};
    \end{tikzpicture}
  \end{center}

  \slidegap
  \aside{Равные ключи \hi{обязаны} давать равные хэши. Ключ --- record,
    record struct или строка: неизменяемые, хэш по содержимому.}
\end{frame}

\begin{frame}{Множество → \texttt{HashSet<T>}}
  \slidetop
  \claim{Словарь \hi{без значений}: тот же хэш, тот же \texttt{Equals},
    ответ --- есть или нет.}

  {\renewcommand{\arraystretch}{1.6}
  \begin{tabular}{@{}p{5.4cm}>{\centering\arraybackslash}p{2.8cm}@{}}
    список в \lo{список} (\texttt{Contains} в цикле) & $O(n^2)$ \\
    список в \hi{множество}                           & $O(n)$ \\
  \end{tabular}}

  \slidegap
  \aside{Объединение, пересечение, разность --- тоже методы, не циклы с
    условиями.}
\end{frame}

\begin{frame}{Отсортированные коллекции}
  \slidetop
  \versus{SortedSet / SortedDictionary}{%
    \begin{cbul}{accent}
      \item внутри дерево
      \item поиск, вставка --- $O(\log n)$
      \item всегда по порядку
    \end{cbul}}{SortedList}{%
    \begin{cbul}{alt}
      \item два массива: ключи, значения
      \item поиск --- $O(\log n)$ бинарным
      \item вставка --- $O(n)$, сдвиг
    \end{cbul}}

  \slidegap
  \aside{Та же логика, что у массива и связного списка: \hi{SortedList} ---
    собрали один раз, читаем много; \hi{SortedDictionary} --- часто
    меняется.}
\end{frame}

\begin{frame}{Стек и очередь на массиве}
  \slidetop
  \begin{center}
    \begin{tikzpicture}[font=\small,
      cell/.style={draw=hairline,line width=0.8pt,fill=codebg,
                   minimum width=1.0cm,minimum height=0.75cm,inner sep=0pt},
      ar/.style={-{Straight Barb[length=1.6mm]},line width=0.7pt}]
      \foreach \i in {0,...,4} \node[cell] (s\i) at (\i*1.0,1.0) {};
      \node[accent,font=\scriptsize,anchor=west] at (s4.east)
        {\texttt{Stack<T>}: push/pop с конца};
      \draw[ar,accent] (s4.north) -- ++(0,0.4);

      \foreach \i in {0,...,4} \node[cell] (q\i) at (\i*1.0,-0.6) {};
      \draw[ar,alt] (q0.north) -- ++(0,0.4)
        node[above,font=\scriptsize,text=alt]{dequeue};
      \draw[ar,alt] ++(4.0,-0.2) -- (q4.north)
        node[above,font=\scriptsize,text=alt]{enqueue};
      \node[alt,font=\scriptsize,anchor=west] at (q4.east)
        {\texttt{Queue<T>}: массив по кругу};
    \end{tikzpicture}
  \end{center}

  \slidegap
  \aside{Оба --- на массиве, не на связном списке: ради \hi{локальности}.}
\end{frame}

\begin{frame}{\texttt{Concurrent*}}
  \slidetop
  \claim{Та же структура, другой договор: доступ из нескольких потоков
    \hi{без внешней блокировки}.}

  \versus{обычная коллекция}{%
    \begin{cbul}{alt}
      \item один писатель за раз
      \item запись из двух потоков --- несогласованное состояние
    \end{cbul}}{\texttt{Concurrent*}}{%
    \begin{cbul}{accent}
      \item \texttt{ConcurrentDictionary}, \texttt{ConcurrentQueue}
      \item атомарные «получить или добавить»
    \end{cbul}}

  \slidegap
  \aside{«Получить или добавить» не гарантирует вызов функции \lo{ровно}
    один раз --- подробно в шестой лекции.}
\end{frame}

\begin{frame}{Лестница интерфейсов}
  \slidetop
  \claim{Принимайте самое \hi{слабое}, возвращайте самое \hi{безопасное}.}

  \begin{center}
    \begin{tikzpicture}[font=\scriptsize,
      st/.style={draw=hairline,line width=0.8pt,fill=codebg,
                 minimum height=0.85cm,minimum width=2.3cm,align=center,
                 inner sep=0.25em},
      ar/.style={-{Straight Barb[length=1.6mm]},muted,line width=0.6pt}]
      \node[st,draw=accent] (e) at (0,0)    {\texttt{IEnumerable<T>}};
      \node[st,draw=accent] (rc) at (2.9,0) {\texttt{IReadOnly-}\\
        \texttt{Collection<T>}};
      \node[st,draw=accent] (rl) at (5.8,0) {\texttt{IReadOnly-}\\
        \texttt{List<T>}};
      \draw[ar] (e) -- (rc); \draw[ar] (rc) -- (rl);

      \node[st,draw=alt] (c) at (2.9,-1.6)  {\texttt{ICollection<T>}};
      \node[st,draw=alt] (l) at (5.8,-1.6)  {\texttt{IList<T>}};
      \draw[ar] (e.south) |- (c.west); \draw[ar] (c) -- (l);

      \node[st,draw=muted,fill=paper] (d) at (8.6,-0.8)
        {\texttt{I(ReadOnly)-}\\\texttt{Dictionary}};

      \node[accent,font=\scriptsize\bfseries,above=0.3em] at (rl.north)
        {только читает};
      \node[alt,font=\scriptsize\bfseries,below=0.3em] at (l.south)
        {можно менять};
    \end{tikzpicture}
  \end{center}

  \slidegap
  \aside{\texttt{IEnumerable} ничего не говорит о стоимости и
    повторяемости --- если проходите дважды, материализуйте один раз.}
\end{frame}

\begin{frame}{Моя рекомендация: иммутабельный стиль}
  \slidetop
  \claim{Коллекция-значение --- по умолчанию \hi{неизменяемая}, а не честное
    слово.}

  \versus{обычный List / Dictionary}{%
    \begin{cbul}{alt}
      \item любой держатель ссылки может изменить
      \item гарантия --- на словах
    \end{cbul}}{Immutable* / IReadOnlyList}{%
    \begin{cbul}{accent}
      \item добавление возвращает \hi{новую} версию
      \item гарантия --- по построению
    \end{cbul}}

  \slidegap
  \aside{Не оптимизация, а защита от ошибок, которые ищутся днями: кто-то
    незаметно изменил структуру, пока вы её читали.}
\end{frame}

\begin{frame}{\texttt{ImmutableList<T>} vs \texttt{ImmutableArray<T>}}
  \slidetop
  {\renewcommand{\arraystretch}{1.6}
  \begin{tabular}{@{}p{4.2cm}>{\centering\arraybackslash}p{3.2cm}%
    >{\centering\arraybackslash}p{3.2cm}@{}}
    & \colhead{accent}{ImmutableList} & \colhead{alt}{ImmutableArray} \\
    внутри       & дерево              & обёртка над массивом \\
    добавление   & \hi{дешёвое}, делит структуру & копирует всё \\
    доступ по индексу & $O(\log n)$    & \hi{O(1)} \\
  \end{tabular}}

  \slidegap
  \aside{Первый --- для коллекций, которые растут по одному элементу.
    Второй --- для тех, что построили один раз и читают по индексу.}
\end{frame}

\begin{frame}{\texttt{Frozen*}}
  \slidetop
  \claim{Собрали один раз --- читаем \hi{миллионы} раз быстрее обычного
    словаря.}

  \slidegap
  {\renewcommand{\arraystretch}{1.5}
  \begin{tabular}{@{}p{3.4cm}p{8.4cm}@{}}
    \textbf{плата}      & построение дороже обычного словаря \\
    \textbf{выгода}      & чтение быстрее: устройство подобрано под
                           конкретный набор ключей \\
    \textbf{где}         & справочники, загруженные при старте: коды
                           стран, таблица маршрутов \\
  \end{tabular}}
\end{frame}

\begin{frame}{\texttt{Span} / \texttt{Memory} --- окно без копирования}
  \slidetop
  \begin{center}
    \begin{tikzpicture}[font=\small,
      cell/.style={draw=hairline,line width=0.7pt,fill=codebg,
                   minimum width=0.9cm,minimum height=0.75cm,inner sep=0pt}]
      \foreach \i in {0,...,6} \node[cell] (c\i) at (\i*0.9,0) {};
      \draw[accent,line width=1.6pt] ([xshift=1pt,yshift=1pt]c2.south west)
        rectangle ([xshift=-1pt,yshift=-1pt]c4.north east);
      \node[accent,font=\scriptsize,below=0.5em] at (c3.south) {Span};
    \end{tikzpicture}
  \end{center}

  \slidegap
  \aside{Та же память, без копирования, только на стеке. \texttt{Memory}
    --- то же окно, но можно хранить в поле и пронести через
    \texttt{await}.}
\end{frame}

\begin{frame}[plain]
  \keyline{Всё в стандартной библиотеке построено поверх массива:\\ каждая
    структура --- выбор, за что платить.}
\end{frame}

% ========================================================= 2. IENUMERABLE ==
\section{2. IEnumerable как контракт}

\begin{frame}[plain]
  \sectionopener{2}{IEnumerable как контракт}{\hi{Обещание}, что перебрать
    можно; когда и сколько раз --- не сказано.}
\end{frame}

\begin{frame}[fragile]{\texttt{foreach}: вы пишете / компилятор делает}
  \slidetop
  \begin{minipage}[t]{0.42\textwidth}
    \colhead{alt}{вы пишете}\par\vspace{0.5em}
\begin{lstlisting}[style=csharp,basicstyle=\ttfamily\small]
foreach (var item in items)
    Console.WriteLine(item);
\end{lstlisting}
  \end{minipage}\hfill
  \begin{minipage}[t]{0.55\textwidth}
    \colhead{accent}{компилятор делает}\par\vspace{0.5em}
\begin{lstlisting}[style=csharp,basicstyle=\ttfamily\fontsize{7.6}{9.0}\selectfont]
using (var e = items.GetEnumerator())
{
    while (e.MoveNext())
    {
        var item = e.Current;
        Console.WriteLine(item);
    }
}
\end{lstlisting}
  \end{minipage}

  \slidegap
  \aside{\texttt{GetEnumerator} --- курсор перед первым элементом.
    \texttt{MoveNext} сдвигает и говорит, есть ли дальше. \texttt{Current}
    --- элемент под курсором. \texttt{using} --- перечислитель может
    держать ресурс.}
\end{frame}

\begin{frame}[fragile]{\texttt{yield return}: описание, а не результат}
  \slidetop
\begin{lstlisting}[style=csharp,basicstyle=\ttfamily\fontsize{7.4}{8.6}\selectfont]
static IEnumerable<int> Evens(int limit)
{
    Console.WriteLine("started");
    for (var i = 0; i <= limit; i++)
        if (i % 2 == 0) yield return i;
    Console.WriteLine("finished");
}
\end{lstlisting}

  \slidegap
\begin{lstlisting}[style=term,basicstyle=\ttfamily\small]
got sequence
started
0 2 4 6 finished
started
0 2 4 6 finished
\end{lstlisting}
  \aside{Вызов \texttt{Evens} не выполнил \hi{ни строки} тела. Каждый
    \texttt{foreach} запускает его с начала.}
\end{frame}

\begin{frame}{Две ловушки отложенного выполнения}
  \slidetop
  \versus{двойной перебор}{%
    \begin{cbul}{alt}
      \item = двойное выполнение
      \item запрос к базе --- два похода
    \end{cbul}}{переменная между переборами}{%
    \begin{cbul}{alt}
      \item лямбда захватила переменную
      \item изменили её --- изменился результат
    \end{cbul}}

  \slidegap
  \hrulethin{0.42}
  \rulegap
  {\large Нужен результат --- материализуйте \hi{один раз}: \texttt{ToList}
    или \texttt{ToArray}.}
\end{frame}

\begin{frame}[plain]
  \keyline{Запрос --- это описание, а не результат:\\ материализуйте, если
    нужен результат ровно один раз.}
\end{frame}

% ================================================================ 3. LINQ ==
\section{3. LINQ-методы}

\begin{frame}[plain]
  \sectionopener{3}{LINQ-методы}{Фильтрация $\rightarrow$ проекция
    $\rightarrow$ агрегация $\rightarrow$ группировка $\rightarrow$
    \hi{соединение}.}
\end{frame}

\begin{frame}[fragile]{Запрос вместо двадцати строк циклами}
  \slidetop
\begin{lstlisting}[style=csharp,basicstyle=\ttfamily\fontsize{7.0}{8.2}\selectfont]
var top = orders
    .Where(o => o.Status == OrderStatus.Paid)
    .GroupBy(o => o.CustomerId)
    .Select(g => new { CustomerId = g.Key, Total = g.Sum(o => o.Amount) })
    .Where(x => x.Total > 100)
    .OrderByDescending(x => x.Total)
    .Take(3)
    .ToList();
// 1: 150, 3: 120
\end{lstlisting}

  \slidegap
  \aside{Взять заказы. Оставить оплаченные. Сгруппировать. Посчитать.
    Отфильтровать. Отсортировать. Взять три. \hi{Материализовать}. Порядок
    слов --- порядок действий.}
\end{frame}

\begin{frame}{Пять семейств операторов}
  \slidetop
  {\renewcommand{\arraystretch}{1.6}
  \begin{tabular}{@{}p{3.0cm}p{2.6cm}p{6.6cm}@{}}
    фильтрация   & \texttt{\hi{Where}}   & оставить подходящие \\
    проекция     & \texttt{\hi{Select}}  & превратить каждый элемент \\
    агрегация    & \texttt{\hi{Sum}}     & свернуть в одно значение \\
    группировка  & \texttt{\hi{GroupBy}} & разложить по ключу \\
    соединение   & \texttt{\hi{Join}}    & сопоставить два источника \\
  \end{tabular}}

  \slidegap
  \aside{\texttt{SelectMany} --- склеить списки списков в одну ленту;
    \texttt{Any} вместо \texttt{Count > 0}; \texttt{Single} --- утверждение,
    что дубликата быть не может.}
\end{frame}

\begin{frame}[fragile]{Запрос --- значение: композиция}
  \slidetop
\begin{lstlisting}[style=csharp,basicstyle=\ttfamily\small]
IEnumerable<Order> query = orders;

if (onlyPaid)
    query = query.Where(o => o.Status == OrderStatus.Paid);

if (minAmount is { } min)
    query = query.Where(o => o.Amount >= min);

var result = query.OrderBy(o => o.Amount).ToList();
\end{lstlisting}

  \slidegap
  \aside{До последней строки \hi{ничего} не вычислено --- \texttt{query} ---
    описание, которое росло. Каждый необязательный фильтр --- две строки,
    независимые от остальных.}
\end{frame}

\begin{frame}[plain]
  \keyline{Запрос говорит, что вы хотите; цикл говорит, как.\\ Переходите
    на «как» только там, где измерили.}
\end{frame}

% =========================================================== 4. IQUERYABLE =
\section{4. IQueryable}

\begin{frame}[plain]
  \sectionopener{4}{IQueryable}{\hi{IEnumerable} выполняет, \hi{IQueryable}
    описывает.}
\end{frame}

\begin{frame}[fragile]{Func или Expression --- в сигнатуре}
  \slidetop
\begin{lstlisting}[style=csharp,basicstyle=\ttfamily\small]
// System.Linq.Enumerable -- runs in memory
Where<T>(this IEnumerable<T> source, Func<T, bool> predicate);

// System.Linq.Queryable -- describes it for a provider
Where<T>(this IQueryable<T> source, Expression<Func<T, bool>> predicate);
\end{lstlisting}

  \slidegap
  \aside{Когда параметр --- \texttt{Expression}, компилятор превращает
    лямбду не в метод, а в \hi{дерево выражений}: объекты, которые можно
    обойти и прочитать, а не выполнить.}
\end{frame}

\begin{frame}{Дерево выражения}
  \slidetop
  \claim{\texttt{o => o.Status == OrderStatus.Paid} --- не код, а
    \hi{данные}.}

  \begin{center}
    \begin{tikzpicture}[font=\small\ttfamily,level distance=1.3cm,
      sibling distance=4.6cm,
      nd/.style={draw=hairline,line width=0.8pt,fill=codebg,
                 minimum height=0.8cm,inner sep=0.3em}]
      \node[nd,draw=accent] {==}
        child { node[nd] {o.Status} }
        child { node[nd] {OrderStatus.Paid} };
    \end{tikzpicture}
  \end{center}

  \slidegap
  \aside{Провайдер обходит дерево в момент перебора и генерирует SQL:
    условие в \texttt{WHERE}, ключ в \texttt{ORDER BY}.}
\end{frame}

\begin{frame}[plain]
  \keyline{IQueryable не выполняет ваш код ---\\ он описывает его для
    кого-то другого.}
\end{frame}

% ===================================================== 5. EXTENSION METHODS=
\section{5. Extension methods}

\begin{frame}[plain]
  \sectionopener{5}{Extension methods}{Статический метод, который выглядит
    как \hi{метод экземпляра}.}
\end{frame}

\begin{frame}[fragile]{\texttt{this} в параметре}
  \slidetop
\begin{lstlisting}[style=csharp,basicstyle=\ttfamily\small]
static class StringExtensions
{
    public static bool IsBlank(this string? value)
        => string.IsNullOrWhiteSpace(value);
}

input.IsBlank();                   // True
StringExtensions.IsBlank(input);   // the same call
nothing.IsBlank();                 // True, no exception on null
\end{lstlisting}

  \slidegap
  \aside{Компилятор не находит метод у типа, ищет среди расширений с
    подходящим первым параметром и переписывает вызов в обычный
    статический. \hi{Нет} доступа к приватным членам; метод экземпляра
    всегда побеждает.}
\end{frame}

\begin{frame}[plain]
  \keyline{Расширение --- функция от своих аргументов и ничего больше.\\
    Если нужно больше --- это класс с зависимостями.}
\end{frame}

\end{document}
